% Laget automatisk av elevark/lag_elevark.py fra opplegg/reaksjonstid.html. Endringer her blir overskrevet.
\documentclass[11pt]{article}
\usepackage{elevark}
\begin{document}
\elevarkhode{Fysikk 1 · Introduksjon til fysikk}{Reaksjonstid}{https://fysikklab.no/opplegg/reaksjonstid.html}{Opplegget på nett}

Hvor raskt reagerer du? Mål reaksjonstiden din på to måter: med en test her på skjermen og med en linjal som faller. Hvor stor er usikkerheten, og gir metodene samme svar?

\section*{Linjalmetoden}

Du trenger en linjal på 30~cm og en partner. Partneren holder linjalen loddrett, med nullpunktet mellom tommelen og pekefingeren din, uten at du rører den. Uten forvarsel slipper partneren linjalen, og du klyper den så fort du kan. Les av hvor langt linjalen falt, i centimeter.

Linjalen faller fritt fra ro, så strekningen er $s = \tfrac{1}{2}gt^{2}$. Da er reaksjonstiden $t = \sqrt{2s/g}$.

\section*{Oppgaver}

\begin{oppgaveliste}
\oppgave \fase{test}{Mål}Gjør minst ti forsøk med hver metode. Skriv reaksjonstiden din på formen $t = (\text{middelverdi} \pm \text{usikkerhet})\,\mathrm{ms}$, med det største avviket fra middelverdien som usikkerhet.
\svar
\oppgave \fase{test}{Regn}Linjalen faller 15~cm før du klyper den. Hva er reaksjonstiden? Hvor langt faller linjalen hvis reaksjonstiden er 0,25~s?
\svar
\oppgave \fase{forklar}{Forklar}Gir de to metodene samme resultat? Nettlesertesten har en systematisk feil: Det tar litt tid fra datamaskinen ber skjermen skifte farge til fargen faktisk vises, og fra du trykker til datamaskinen registrerer det. Blir målt tid for stor eller for liten?
\svar
\oppgave \fase{forklar}{Vurder}Hvilke tilfeldige og systematiske feil er det i linjalmetoden? Hvorfor er det viktig at partneren ikke gir forvarsel?
\svar
\oppgave \fase{test}{Regn}Du kjører bil i 80~km/h. Hvor langt kjører du i løpet av reaksjonstiden din, før du begynner å bremse? Gjør om til SI-enheter først.
\svar
\oppgave \fase{utfordring}{Utfordring}Linjalen er 30~cm lang. Hva er den lengste reaksjonstiden linjalmetoden kan måle? Hvordan ville du laget en «reaksjonslinjal» med tider i stedet for centimeter? Er avstanden mellom 50 ms-strekene den samme hele veien?
\svar
\end{oppgaveliste}

\end{document}
